% !TeX Program = lualatex

\documentclass[a4paper]{article}
\usepackage[left=2cm,right=2cm,top=1cm,bottom=1cm]{geometry}
\usepackage{fontspec}
\usepackage{libertinus}
\renewcommand*\familydefault{\sfdefault}

\usepackage{xcolor}
\colorlet{textgreen}{black!50!green!50!white}
\colorlet{bggreen}{black!50!green!10!white}
\colorlet{xgreen}{black!50!green!10!white}
\colorlet{linegreen}{black!50!green}

\usepackage{pst-barcode}

\usepackage{tikz}
\usetikzlibrary{calc}
\usetikzlibrary{shapes.geometric}

\usepackage{enumitem}

\usepackage{tabularx}
\usepackage{colortbl}
\usepackage{booktabs}
\usepackage{array}
\arrayrulecolor{linegreen}
\setlength{\aboverulesep}{0pt}
\setlength{\belowrulesep}{0pt}
\newcommand{\PreserveBackslash}[1]{\let\temp=\\#1\let\\=\temp}
\newcolumntype{C}[1]{>{\PreserveBackslash\centering}p{#1}}
\newcolumntype{R}[1]{>{\PreserveBackslash\raggedleft}p{#1}}
\newcolumntype{L}[1]{>{\PreserveBackslash\raggedright}p{#1}}
\newcolumntype{Y}{>{\centering\arraybackslash}X}

\usepackage{calc}
\usepackage{everypage}
\usepackage{lastpage}
\usepackage{multicol}

\def\documentid{P-A38-2k25}

\AddEverypageHook{%
    \begin{tikzpicture}[remember picture, overlay]
        \node [color=textgreen, rotate=90, anchor=north] at ($(current page.west)+(10mm,0)$)
        {\ttfamily\bfseries Originalpapier nur, wenn dieser Hinweis in Gründruck erscheint.};
        \node [rotate=90, anchor=north west, align=left] at ($(current page.south west)+(10mm,10mm)$)
        {%
            \hspace{1cm}\texttt{\documentid}
            \\
            \hspace{1cm}\texttt{Seite \thepage{} von \pageref{LastPage}}
            \\
            \psbarcode[transy=0.35,scalex=0.5,scaley=0.5]{\documentid-\thepage}{parse,format=square: version=18x18}{datamatrix}
        };
    \end{tikzpicture}%
    %
    \def\foldFirst{105mm}%
    \def\foldSecond{210mm}%
    \begin{tikzpicture}[remember picture, overlay]
        \node (foldAA) at ($(current page.north west) + (5mm, - \foldFirst)$) {};
        \node (foldAB) at ($(current page.north west) + (10mm, - \foldFirst)$) {};
        \node (foldBA) at ($(current page.north west) + (5mm, - \foldSecond)$) {};
        \node (foldBB) at ($(current page.north west) + (10mm, - \foldSecond)$) {};
        \node (centerA) at ($(current page.west) + (5mm, 0)$) {};
        \node (centerB) at ($(current page.west) + (10mm, 0)$) {};
        \draw[-] (foldAA) -- (foldAB);
        \draw[-] (foldBA) -- (foldBB);
        \draw[-] (centerA) -- (centerB);
    \end{tikzpicture}%
}

\parindent=0pt
\pagestyle{empty}

\usepackage[ngerman]{babel}

\newcommand\sectitle[1]{\textbf{\Large\raisebox{0pt}[1.5\height][1.5\depth]{\vphantom{#1}}#1}}
\newcommand\tftitle[1]{\textbf{#1}}
\newcommand\tfnote[1]{\textit{#1}}
\newcommand\X{\begin{tabular}{|>{\columncolor{white}}l|}\hline\raisebox{0pt}[1.5\height][1.5\depth]{\vphantom{Fg}\hspace{-0.29em}\color{xgreen}X\hspace{-0.29em}}\\\hline\end{tabular}}
\newcommand\digittf{\begin{tabular}{|>{\columncolor{white}}l|}\hline\raisebox{0pt}[1.5\height][1.5\depth]{\vphantom{Fg}}\\\hline\end{tabular}}
\newcommand\widetf{\begin{tabular}{|>{\columncolor{white}}l|}\hline\raisebox{0pt}[1.5\height][1.5\depth]{\vphantom{Fg}}\hspace*{\linewidth - 1.5em}\\\hline\end{tabular}}
\newcommand\Xsub[1]{$\underset{#1}{\X}$}

\newcounter{ansnum}
\newcommand\ans{\refstepcounter{ansnum}\arabic{ansnum}}

\begin{document}%
\begin{tikzpicture}[remember picture, overlay]
    \node [rotate=90, anchor=south] at ($(current page.east)+(-15mm,6.75cm)$)
    {\ttfamily\bfseries\footnotesize Die kleingedruckten Randbemerkungen bitte nicht lesen.};
    \node [anchor=south west] at ($(current page.south west)+(0,5mm)$)
    {%
        \hspace*{1cm + 0.33ex}%
        \ttfamily%
        blinry.org/passierschein-a38
        \textrightarrow{}
        hans.coffee
        \cdot{}
        CC BY-SA 4.0
        \cdot{}
        Weg zum Hafen bitte umseitig zur Kenntnis nehmen
    };
\end{tikzpicture}%
Antragsformular für den

\begin{Huge}
    Passierschein A38
\end{Huge}

~

\begin{tabularx}{\linewidth}{@{}Xr@{}}
    \tftitle{Laufende Nummer} \tfnote{(Bitte schätzen)}                                              & nach Rundschreiben B64            \\
    \digittf\digittf\digittf\digittf\digittf\digittf\digittf\digittf\digittf\digittf\digittf\digittf & der Bundesbehörde für Rautavistik
\end{tabularx}

\vspace*{3ex}

\hspace*{-1cm}\begin{tabularx}{\linewidth + 1cm}{@{} R{7.5mm} @{\hspace*{2.5mm}}| >{\columncolor{bggreen}}X |@{}}
    \cmidrule{2-2}
         & \sectitle{Persönliche Angaben}                                                                \\
    \ans & \tftitle{Was ist Ihrer Ansicht nach die ideale Anzahl von Schafen?}                           \\
         & \tfnote{Bitte beachten Sie, dass Ziegen grundsätzlich keine Schafe sind.}                     \\
         & \digittf\digittf\digittf\digittf                                                              \\
         &                                                                                               \\
    \ans & \tftitle{Beschreiben Sie -- in vier Wörtern oder weniger -- Ihre Meinung zu Plattentektonik.} \\
         & \widetf                                                                                       \\
         &                                                                                               \\
    \ans & \tftitle{Nennen Sie eine Ihrem Charakter entsprechende Salatsorte.}                           \\
         & \widetf                                                                                       \\
         &                                                                                               \\
    \ans & \tftitle{Können Sie fliegen?}                                                                 \\
         & \X{} Ja \qquad \X{} Momentan nicht.                                                           \\
         &                                                                                               \\
    \ans & \tftitle{Wie flexibel sind Ihre Zehen?}                                                       \\
         & Mäßig flexibel \Xsub{1} \Xsub{2} \Xsub{3} \Xsub{4} \Xsub{5} Sehr flexibel                     \\
         &                                                                                               \\
    \cmidrule{2-2}
         & \sectitle{Gesetzliche Vorschriften}                                                           \\
    \ans & \tftitle{Führen Sie Waffen bei sich?}                                                         \\
         & \X{} Schere \qquad \X{} Stein \qquad \X{} Papier                                              \\
         &                                                                                               \\
    \ans & \tftitle{Summen Sie die ersten vier Noten Ihres Lieblingsliedes.}                             \\
         & \X{} Mhmmm-mmhhh hmm hmmmmmm                                                                  \\
         & \X{} Wupdup dup dup                                                                           \\
         & \X{} Drrrrnngggggggg krtz-krtz-krtz                                                           \\
         & \X{} Aaaahhh na-na-naaaaaa                                                                    \\
         &                                                                                               \\
    \ans & \tftitle{Bitte kreuzen Sie dieses Feld an:}                                                   \\
         & \X{} Nein!                                                                                    \\
         &                                                                                               \\
    \ans & \tftitle{Haben Sie in den letzten 14 Monaten:}                                                \\
         & \tfnote{Nur Einfachantwort möglich.}                                                          \\
         & \X{} bereits eine Quellenbeihilfeerklärung abgegeben                                          \\
         & \X{} doppelte Substanzkennzahlen in Anspruch genommen                                         \\
         & \X{} (auch zeitweise) die Interessensverordnungen gelesen                                     \\
         & \X{} einen Passierschein A38 beantragt                                                        \\
         &                                                                                               \\
    \ans & \tftitle{Die Anzahl der angekreuzten Felder auf diesem Formular ist:}                         \\
         & \X{} gerade                                                                                   \\
         & \X{} ungerade                                                                                 \\
         &                                                                                               \\
    \cmidrule{2-2}
\end{tabularx}

\begin{multicols}{3}
    \tftitle{Hinweise}\\
    \begin{footnotesize}
        Dieser Antrag wird maschinell eingelesen und ist auch ohne Ihre Unterschrift gültig.

        Mit Berühren dieses Formulars stimmen Sie der Verarbeitung Ihrer DNA zu.

        Eine Braille-Version dieses Formulars senden wir Ihnen auf Wunsch gern als PDF zu.
    \end{footnotesize}

    \vspace*{2.5ex}\tftitle{Freifeld für zusätzliche Entropie:}
    \begin{tabularx}{\linewidth}{|Y|}
        \midrule
        \\
        \tfnote{(\ans{} Bei Bedarf bekritzeln.)} \\
        \\
        \midrule
    \end{tabularx}

    \columnbreak

    \tftitle{Bitte fügen Sie folgende Anlagen bei:}
    \begin{itemize}[noitemsep,topsep=0pt,parsep=0pt,partopsep=0pt]
        \item Blaues Formular von Schalter I.
        \item Rosa Formular von Schalter XII, Stiege B, Korridor J.
        \item Ausgefülltes Antragsformular für den Passierschein A38.
    \end{itemize}

    \begin{tabularx}{\linewidth}{@{}lX@{}}
        \ans{} \tftitle{Nasenabdruck}
         &
        \tikz \draw[dashed,linegreen]
        (0,0) to[out=0, in=-90]
        (0.66,0.33) to[out=90, in=0]
        (0,1.5) to[out=180,in=90]
        (-0.66,0.33) to[out=-90,in=180]
        (0,0);
    \end{tabularx}

    \columnbreak

    \begin{itemize}[noitemsep]
        \item[\ans{} \X] Ich verzichte ausdrücklich auf die Entbindung von der Ausfüllungspflicht.
        \item[\ans{} \X] Ich ziehe meinen Antrag zurück. \tfnote{(Pflichtfeld)}
        \item[\ans{} \X] Batman-Regelung?
            \\[4ex]
        \item[\ans{} \phantom{\X}] \dotfill\\
            Unterwäsche
            \smash{
                \begin{tikzpicture}[remember picture, overlay]
                    \node [
                        draw,
                        dashed,
                        linegreen,
                        ellipse,
                        align=center,
                        minimum width=6cm,
                        minimum height=3cm,
                        inner sep=0pt,
                        rotate=90
                    ]
                    at ($(current page.east)+(5mm,-11cm)$) {
                        \color{black} \ans{} \\
                        \color{black} Hier abbeißen.
                        \\
                        \\
                        \\
                        \\
                        \\
                        \\
                    };
                \end{tikzpicture}
            }
    \end{itemize}%
\end{multicols}%
\clearpage%

Der Weg zum Hafen kann standortabhängig unterschiedlich sein.
In der Regel befindet sich der Hafen unterhalb der Stadt -- immer da wo Wasser ist.
Beachten Sie die örtlichen Aushänge oder fragen Sie das Hafenpersonal.

\vfill
\begin{center}
    \tfnote{(Diese Seite wurde unabsichtlich freigelassen)}
\end{center}
\vfill
\vfill
\vfill
\vfill

\end{document}